% ==============================================================================
% Fichier : front/02_mode_emploi.tex : FORN (LPI342)
% ==============================================================================

\chapter*{Mode d'emploi}
\addcontentsline{toc}{chapter}{Mode d'emploi}
\markboth{Mode d'emploi}{Mode d'emploi}

\section*{À qui s'adresse ce cours}

Aux étudiants du Master Professionnel en Sécurité des Systèmes Informatiques
et des Données, et aux praticiens qui exercent ou se destinent aux métiers
d'\textbf{analyste forensique}, d'\textbf{expert judiciaire},
d'\textbf{enquêteur en cybercriminalité} et de \textbf{juriste TIC}. Des
notions de réseaux, de systèmes d'exploitation et de droit numérique sont
utiles.

\section*{Comment lire ce livre}

Le livre comprend un prologue et six parties. Chaque chapitre suit le même
plan~: épigraphe, niveau et prérequis, objectifs, question directrice, sections,
« à retenir », dix exercices gradués, transition. Les niveaux des exercices
sont signalés par des étoiles~: \niveaubase{} fondamental, \niveaumoyen{}
intermédiaire, \niveauexpert{} expert.

\section*{Les encadrés}

\begin{itemize}
    \item \textbf{Scénario opérationnel}~: un épisode de l'opération TSANAGA,
          avec \emph{constatation} (faits) et \emph{analyse} (interprétation,
          niveau de confiance).
    \item \textbf{Artefact}~: la fiche d'une pièce (nature, source, pertinence,
          fiabilité, opposabilité).
    \item \textbf{Chaîne de custody}~: une procédure de préservation de la
          preuve.
    \item \textbf{Analyse CRO}~: l'arbitrage Confidentialité, Fiabilité,
          Opposabilité.
    \item \textbf{Application métier}~: ce que le chapitre change pour un
          métier.
    \item \textbf{Cadre juridique}~: le texte de droit applicable.
\end{itemize}

\section*{Constatation et analyse}

Le cours distingue systématiquement ce que l'on \textbf{constate} (un fait
vérifiable) de ce que l'on en \textbf{déduit} (une interprétation, assortie
d'un niveau de confiance). C'est la discipline première de l'enquête.

\section*{NDAP SUKULU et évaluation}

L'environnement compagnon NDAP SUKULU propose, pour chaque chapitre, une fiche
d'auto-évaluation et, pour les chapitres pratiques, un lab sur pièces. Fiches
et labs se font hors des 30~heures de cours. Le code du cours est
\textbf{\codeUE}.
